\section{Related Work}
\label{sec:related}

\paragraph{Linguistic cues for fake news.}
A long line of work characterises deceptive or unreliable news through stylistic and psycholinguistic
signals. \citet{horne2017fake} show that fake articles have longer, noun-heavy titles and simpler,
more repetitive bodies; \citet{rashkin2017truth} find that subjective, hedging and superlative
language separates unreliable from trusted sources; \citet{perezrosas2018automatic} combine
$n$-grams, readability, syntax and LIWC categories. \citet{wang2017liar} release LIAR and combine
text with speaker metadata. Surveys by \citet{shu2017fake} and \citet{zhou2020survey} group these
signals under style-based detection and note their appeal for explainability, but also their fragility
across domains. Almost all of this work targets English, whose relevant markers (modal verbs, hedges,
function-word ratios) do not transfer directly to an agglutinative language such as Indonesian.

\paragraph{Indonesian hoax detection.}
Early Indonesian systems rely on bag-of-words classifiers; \citet{pratiwi2017hoax}, for instance,
apply naïve Bayes to a small corpus of referee-labelled news, which we reuse as an out-of-distribution
test set. Pre-trained Indonesian encoders such as IndoBERT \citep{wilie2020indonlu,koto2020indolem}
have since become the default for Indonesian text classification. Neither line of work models
Indonesian morphology or evidential marking explicitly, and results are typically reported on a single
in-domain split, which leaves open whether high scores reflect transferable cues or corpus artefacts.

\paragraph{Knowledge-informed fusion.}
Injecting hand-crafted features into neural text classifiers is commonly done by concatenation.
Gated fusion, popularised for multimodal learning by the gated multimodal unit
\citep{arevalo2017gated}, learns input-dependent mixing weights instead. \menet{} applies a
complementary gate between a learned contextual view and a symbolic linguistic view of the same text.

\paragraph{Shortcut learning and calibration.}
Neural classifiers readily exploit dataset artefacts that correlate with labels
\citep{gururangan2018annotation,mccoy2019right,geirhos2020shortcut}. For misinformation corpora built
by merging sources (e.g.\ fact-checking sites for hoaxes and mainstream outlets for valid news), length,
formatting and source boilerplate are obvious candidates. We therefore evaluate with a format-normalised
control and a cross-dataset test, and report calibration \citep{naeini2015ece,guo2017calibration} in
addition to accuracy, since a moderation tool that is confidently wrong is worse than one that abstains.
